\documentclass[11pt]{article}
\usepackage[T1]{fontenc}
\usepackage[utf8]{inputenc}
\usepackage{lmodern}
\usepackage[margin=0.95in]{geometry}
\usepackage{amsmath,amssymb,amsthm}
\usepackage{booktabs}
\usepackage{microtype}
\usepackage[colorlinks=true,linkcolor=blue,citecolor=blue,urlcolor=blue]{hyperref}

\allowdisplaybreaks

\newtheorem{theorem}{Theorem}[section]
\newtheorem{proposition}[theorem]{Proposition}
\newtheorem{lemma}[theorem]{Lemma}
\theoremstyle{remark}
\newtheorem{remark}[theorem]{Remark}

\title{A Sharp Local Bound for the Fifth Complete Pivot}
\author{Claude Make No Mistakes}
\date{November 2026}

\begin{document}
\maketitle

\begin{abstract}
We study a one-tie relaxation of the extremal complete-pivoting pattern of
Chen, Edelman, and Urschel for real $5\times 5$ matrices, retaining all
fourteen prescribed initial $\pm 1$ entries and all stage-2, stage-3, and
off-diagonal stage-4 ties while relaxing the final diagonal tie $H_{22}=H_{11}$
to $-1\le r = H_{22}/H_{11} < 1$. Over an explicit rational parameter box
$B$ surrounding the published extremal matrix, we prove via exact rational
interval arithmetic, symbolic elimination, and exact polynomial root
certification that every matrix in the relaxed family has fifth pivot strictly
below the degree-61 benchmark $g_* \approx 4.13251707863247$, that the
supremum over the relaxed family in $B$ equals $g_*$, and that $g_*$ is
attained in $B$ when $r=1$ is permitted. We also certify an explicit algebraic
witness in $B$ with $r \approx 0.999939320464152$ and
$p_5 \approx 4.1324652818749806$, and we establish a global necessary
inequality over the entire relaxed family: any counterexample with $|p_5|>g_*$
must satisfy $0<z<1$ and $(1+r)(2+z(1-z))>g_*$, hence
$r > 4g_*/9 - 1 \approx 0.83667$.
\end{abstract}

\section{Introduction and Main Results}\label{sec:intro}

In Gaussian elimination with complete pivoting on a real $5\times 5$ matrix
$A$, one normalizes the initial matrix so that $A_{11}=1$ and
$|A_{ij}|\le 1$ for all $i,j$, and requires at each elimination stage $k\in\{1,2,3,4\}$
that the pivot $p_k = A^{(k)}_{11}$ is nonzero and dominates every entry of
the remaining $(6-k)\times(6-k)$ Schur complement:
\[
|A^{(k)}_{ij}| \le |A^{(k)}_{11}| \qquad (1\le i,j\le 6-k).
\]
Chen, Edelman, and Urschel~\cite{CEU2026} introduced a structured equality
pattern (their equation~(2.1)) and conjectured that the fifth pivot of every
completely pivoted real $5\times 5$ matrix is bounded by a degree-61 algebraic
number $g_* \approx 4.13251707863247$ (the unique root in $[4,5]$ of the
integer polynomial $P_5$ in~\cite[eq.~(2.15)]{CEU2026}).

Write the successive Schur complements as $F = A^{(2)}\in\mathbb{R}^{4\times 4}$,
$G = A^{(3)}\in\mathbb{R}^{3\times 3}$, and $H = A^{(4)}\in\mathbb{R}^{2\times 2}$,
indexing each stage from $1$, and let $p_5 = A^{(5)}_{11}$ denote the fifth
pivot. We study the following one-tie relaxation of the pattern
in~\cite[eq.~(2.1)]{CEU2026}:
\begin{enumerate}
  \item \emph{Fourteen prescribed initial entries:}
  \begin{equation}\label{eq:template}
  A = \begin{pmatrix}
    1 & 1 & a_{13} & a_{14} & a_{15} \\
    x & a_{22} & a_{23} & -1 & 1 \\
    -1 & a_{32} & a_{33} & -1 & -1 \\
    y & a_{42} & 1 & 1 & 1 \\
    z & -1 & 1 & -1 & 1
  \end{pmatrix},
  \end{equation}
  where we designate $x=A_{21}$, $y=A_{41}$, and $z=A_{51}$ as the free
  parameters.
  \item \emph{Nine retained Schur-complement ties:}
  \begin{align*}
    \text{Stage 2:}\quad & F_{21}=F_{11},\quad F_{23}=-F_{11},\quad
                           F_{32}=F_{11},\quad F_{41}=-F_{11}, \\
    \text{Stage 3:}\quad & G_{13}=-G_{11},\quad G_{21}=G_{11},\quad
                           G_{32}=-G_{11}, \\
    \text{Stage 4:}\quad & H_{12}=H_{11},\quad H_{21}=-H_{11}.
  \end{align*}
  \item \emph{Relaxed stage-4 diagonal ratio:} instead of imposing
  $H_{22}=H_{11}$, we require
  \[
    -1 \le r := \frac{H_{22}}{H_{11}} < 1.
  \]
\end{enumerate}
Let $\mathcal{F}_{<}$ denote the set of all real completely pivoted $5\times 5$
matrices with nonzero pivots $p_1,\dots,p_5$ satisfying conditions~(1)--(3),
and let $\mathcal{F}_{\le}$ denote the corresponding family when the boundary
tie $r=1$ ($H_{22}=H_{11}$) is also permitted.

Throughout, every terminating decimal denotes an exact rational number. Define
the closed rational parameter box $B = [x_L,x_R]\times[y_L,y_R]\times[z_L,z_R]$
given in Table~\ref{tab:box}.

\begin{table}[htbp]
\centering
\caption{The rational parameter box $B$ for $(x,y,z)=(A_{21},A_{41},A_{51})$.}
\label{tab:box}
\begin{tabular}{@{}ll@{}}
\toprule
Parameter & Closed interval \\
\midrule
$x = A_{21}$ & $[x_L, x_R] = [-0.6177,\, -0.6174]$ \\
$y = A_{41}$ & $[y_L, y_R] = [-0.7794,\, -0.7789]$ \\
$z = A_{51}$ & $[z_L, z_R] = [0.45322490,\, 0.45322501]$ \\
\bottomrule
\end{tabular}
\end{table}

Our main results are the following three statements.

\begin{theorem}[Sharp local bound and attainment on $B$]\label{thm:box}
Let $g_*$ be the unique real root in $[4,5]$ of the degree-61 benchmark
polynomial $P_5$ of~\cite[eq.~(2.15)]{CEU2026}.
\begin{enumerate}
  \item Every matrix $A\in\mathcal{F}_{<}$ with $(A_{21},A_{41},A_{51})\in B$
  satisfies $0 < p_5(A) < g_*$.
  \item The local supremum is sharp:
  \[
    \sup\bigl\{p_5(A) : A\in\mathcal{F}_{<},\; (A_{21},A_{41},A_{51})\in B\bigr\} = g_*.
  \]
  \item There exists a full-equality matrix $A_*\in\mathcal{F}_{\le}$ with
  $(A_{21},A_{41},A_{51})\in B$ and $r=1$ that attains $p_5(A_*)=g_*$.
\end{enumerate}
\end{theorem}

\begin{theorem}[Global necessary condition on $\mathcal{F}_{<}$]\label{thm:global}
Every matrix $A\in\mathcal{F}_{<}$ satisfies $p_5(A)>0$. If $A\in\mathcal{F}_{<}$
satisfies $p_5(A) > g_*$, then $0 < z < 1$ and
\begin{equation}\label{eq:necessary}
  (1+r)\bigl(2+z(1-z)\bigr) > g_*.
\end{equation}
In particular, any counterexample in $\mathcal{F}_{<}$ must satisfy
\[
  r > \frac{4g_*}{9} - 1 \approx 0.83667
  \qquad\text{and}\qquad
  z(1-z) > \frac{g_*}{2} - 2 \approx 0.06626.
\]
\end{theorem}

\begin{proposition}[Certified nonemptiness witness in $B$]\label{prop:witness}
There exists an explicit matrix $A\in\mathcal{F}_{<}$ with
$x = -386/625$, $z = 18129/40000$, and $y$ equal to the unique root in
$(-0.779073,-0.779072)$ of the irreducible cubic polynomial~\eqref{eq:witness-cubic},
whose stage-4 ratio and fifth pivot satisfy
\begin{align*}
  0.99993932046415197561 &< r < 0.99993932046415197562, \\
  4.13246528187498059056 &< p_5 < 4.13246528187498059057.
\end{align*}
\end{proposition}

The region of $\mathcal{F}_{<}$ outside the box $B$ that satisfies the
necessary condition~\eqref{eq:necessary} remains open in this work.

\section{Benchmark Isolation and Relaxed Witness}\label{sec:benchmark}

The 62 integer coefficients of the degree-61 benchmark polynomial $P_5$
archived in \path{pivot/p5_coefficients.txt} agree term by term with
\cite[eq.~(2.15)]{CEU2026} and \texttt{Out[12]} of the authors' discriminant
notebook~\cite{CEUCode}. Given an integer polynomial $f$ of degree $n$ and
rationals $\alpha < \beta$, the M\"obius substitution
$u = (\alpha + \beta t)/(1+t)$ maps $(0,\infty)$ bijectively onto
$(\alpha,\beta)$ while preserving root multiplicities; clearing denominators
yields an integer polynomial whose sign variations bound the root count on
$(\alpha,\beta)$ by Descartes' rule of signs. Applying this exact substitution
to $P_5$ certifies:
\begin{itemize}
  \item $1$ sign variation and opposite endpoint signs on $[4,5]$, hence a
  unique root $g_*\in(4,5)$;
  \item $0$ sign variations and nonzero endpoints on $[-5,0]$, hence no root
  in $[-5,0]$;
  \item $1$ sign variation and opposite endpoint signs on $[g_L,g_H]$, giving
  the strict 30-digit rational bracket
  \begin{equation}\label{eq:gstar-bracket}
  \begin{aligned}
    g_L &:= 4.132517078632472854223346853277 < g_* \\
        &< 4.132517078632472854223346853278 =: g_H.
  \end{aligned}
  \end{equation}
\end{itemize}

As a verification control, the exact full-equality matrix at
$x=y=-1/2$, $z=0$ has pivots $1,1,2,2,4$ and $r=1$. The exact checker
(\texttt{pivot/verify.py}) accepts this control when $r=1$ is permitted,
rejects it under the strict $r<1$ criterion, and rejects perturbed controls
with a wrong prescribed entry or a stage-1 dominance violation.

\begin{proof}[Proof of Proposition~\ref{prop:witness}]
Set $x=-386/625=-0.6176$ and $z=18129/40000=0.453225$. Substituting into the
parametric reconstruction of Section~\ref{sec:box-proof} and clearing
denominators in the equation $H_{11}+H_{21}=0$ produces the primitive
irreducible cubic polynomial
\begin{equation}\label{eq:witness-cubic}
\begin{aligned}
  C(y) &= 3586674629163499482445312500000000000\,y^3 \\
       &\quad - 22693677055802856449945261748450390625\,y^2 \\
       &\quad - 21794053267721803248556940294525366975\,y \\
       &\quad - 1509132804792772878470239805058936768.
\end{aligned}
\end{equation}
Exact Sturm root counting certifies that $C$ is irreducible over $\mathbb{Q}$,
is nonzero at $-0.779073$ and $-0.779072$, and has a unique real root
$y_w \in (-0.779073,-0.779072)$.

Working in the cubic number field $\mathbb{Q}[y]/(C(y))$ with exact rational
interval refinement around $y_w$, the certificate (\texttt{pivot/verify.py},
archived in \texttt{pivot/witness.json}) verifies that every denominator in
the reconstruction is invertible modulo $C$, checks every entry of all five
Schur stages against its independent bordered-determinant ratio from $A$,
verifies all fourteen prescribed entries and all nine retained ties to be
identically zero in $\mathbb{Q}[y]/(C(y))$, proves strict positivity of all
five pivots and nonnegativity of both dominance slacks
$|p_k|\pm A^{(k)}_{ij}\ge 0$ at every entry, and certifies $-1\le r<1$
together with the stated 20-digit enclosures for $r$ and $p_5$.
\end{proof}

As an independent local check requiring no elimination or discriminant
computation, direct rational interval arithmetic
(\path{verification/pivot_subbox.py}) on the witness sub-box
\[
  [-0.6176001,\,-0.6175999]
  \times [-0.779073,\,-0.779072]
  \times [0.45322499,\,0.45322501]
  \subset B
\]
certifies directly from the Schur formulas that every feasible point in that
sub-box satisfies $4.132464 < p_5 < 4.132467 < g_L < g_*$.

\section{Proof of the Local Upper Bound on the Box}\label{sec:box-proof}

We now prove Part~1 of Theorem~\ref{thm:box}. Define the rational functions of
$(x,y,z)$:
\begin{equation}\label{eq:aux-vars}
\begin{gathered}
  a = 1+x,\qquad b = 1+z,\qquad q = 2-z(x+z),\\
  t = -\frac{z}{y},\qquad v = \frac{z(x+z)}{a},\qquad d = -1+xz,\qquad e = 1-xv,
\end{gathered}
\end{equation}
and set
\[
  k = \frac{q-za-y(z+v)}{e-d},\qquad u = \frac{b-q}{k}.
\]

\begin{lemma}[Reconstruction on $B$]\label{lem:reconstruct}
Throughout the box $B$, the rational functions $a,b,e-d,k,q,H_{11}$ are
strictly positive and $y$ is strictly negative. Every matrix
$A\in\mathcal{F}_{\le}$ with $(A_{21},A_{41},A_{51})=(x,y,z)\in B$ is uniquely
determined by $(x,y,z)$ as
\begin{equation}\label{eq:reconstructed-A}
A = \begin{pmatrix}
  1 & 1 & t & -z & v \\
  x & a+z & u+xt & -1 & 1 \\
  -1 & z & q+u-t & -1 & -1 \\
  y & bk+y & 1 & 1 & 1 \\
  z & -1 & 1 & -1 & 1
\end{pmatrix}.
\end{equation}
Conversely, for every $(x,y,z)\in B$, the matrix~\eqref{eq:reconstructed-A}
has well-defined Schur complements $F,G,H$, satisfies all fourteen prescribed
entries, all stage-2 and stage-3 ties, the stage-4 tie $H_{12}=H_{11}$, and
strict pivot dominance at all non-tied entries of stages~$1,2,3$.
\end{lemma}

\begin{proof}
Let $A\in\mathcal{F}_{\le}$ have $(A_{21},A_{41},A_{51})=(x,y,z)\in B$.
Directly from $F_{ij}=A_{i+1,j+1}-A_{i+1,1}A_{1,j+1}$ and~\eqref{eq:template},
$F_{41}=-1-z=-b$. The four stage-2 ties give:
\begin{itemize}
  \item $F_{41}=-F_{11} \implies F_{11}=b$, so $a_{22}=F_{11}+x=a+z$;
  \item $F_{21}=F_{11} \implies a_{32}+1=1+z$, so $a_{32}=z$;
  \item $F_{23}=-F_{11} \implies -1+a_{14}=-(1+z)$, so $a_{14}=-z$ and
  hence $F_{13}=-1-xa_{14}=d$;
  \item $F_{32}=F_{11} \implies 1-ya_{13}=1+z$, so $a_{13}=-z/y=t$ (since $y<0$ on $B$).
\end{itemize}
At stage~3, using $F_{11}=b>0$ and $F_{21}=b$, $F_{41}=-b$, we have
$G_{32}=F_{43}+F_{13}=(-1+z^2)+(-1+xz)=-2+z(x+z)$, so the tie
$G_{32}=-G_{11}$ forces $G_{11}=q$. Next,
$G_{13}=F_{24}-F_{14}=-2+(1+x)a_{15}$, so the tie $G_{13}=-q$ gives
$a_{15}=z(x+z)/a=v$ (since $a>0$ on $B$) and hence $F_{14}=1-xv=e$.
Writing $k:=F_{31}/b$, the ties $G_{21}=q$ and $G_{11}=q$ give
$kF_{12}=b-q$ and $F_{22}=q+F_{12}$. Finally, at stage~4,
$G_{12}=-b-d=-za$, so the retained tie $H_{12}=H_{11}$ reads
\[
  1 - yv - ke + q = 1 + yz - kd + za,
\]
which determines $k=(q-za-y(z+v))/(e-d)$ and $F_{12}=u=(b-q)/k$, yielding the
entries $a_{23}=u+xt$, $a_{33}=q+u-t$, and $a_{42}=bk+y$ of~\eqref{eq:reconstructed-A}.

In the forward direction, exact symbolic substitution (\texttt{pivot/region.py})
verifies that~\eqref{eq:reconstructed-A} produces the compact Schur
complements
\begin{align*}
F &= \begin{pmatrix}
  b & u & d & e \\
  b & q+u & -b & -1+v \\
  bk & b & 1+yz & 1-yv \\
  -b & 1-zt & -1+z^2 & 1-zv
\end{pmatrix}, &
G &= \begin{pmatrix}
  q & -za & -q \\
  q & 1+yz-kd & 1-yv-ke \\
  1-zt+u & -q & 2-(x+z)v
\end{pmatrix},
\end{align*}
and
\[
H = \begin{pmatrix}
  G_{22}-G_{12} & G_{23}+q \\
  -q - G_{31}G_{12}/q & G_{33}+G_{31}
\end{pmatrix},
\]
and satisfies all stage-2 and stage-3 ties together with $H_{12}=H_{11}$.
Exact rational interval arithmetic (\texttt{pivot/region.py}, archived in
\texttt{pivot/region.json}) proves $a,b,e-d,k,q,H_{11}>0$ and $y<0$ throughout
$B$, so every rational expression is $C^\infty$ on an open neighborhood of $B$.
Applying the mean-value interval enclosure to every non-identically-zero slack
$A^{(k)}_{11}\pm A^{(k)}_{ij}$ for $k=1,2,3$ certifies that all stage-1
through stage-3 dominance inequalities hold on all of $B$.
\end{proof}

Define on $B$ the three smooth functions
\[
  P := H_{11}+H_{21},\qquad D := H_{22}-H_{11},\qquad W := H_{11}+H_{22}.
\]
Because $H_{12}=H_{11}$ identically on $B$, whenever $P=0$ (i.e., $H_{21}=-H_{11}$)
the fifth pivot is
\[
  p_5 = H_{22} - \frac{H_{21}H_{12}}{H_{11}} = H_{22} + H_{11} = W.
\]

\begin{lemma}[Unique $P=0$ branch and strict comparator reduction]\label{lem:branch}
For each fixed $z\in[z_L,z_R]$:
\begin{enumerate}
  \item There is a unique function $y(\cdot,z)\colon [x_L,x_R]\to (y_L,y_R)$
  such that $P(x,y(x,z),z)=0$. It is $C^\infty$ on $[x_L,x_R]\times[z_L,z_R]$
  with $\partial y/\partial x = -P_x/P_y$.
  \item Both $\tilde{D}(x):=D(x,y(x,z),z)$ and $\tilde{W}(x):=W(x,y(x,z),z)$
  are strictly increasing in $x$ on $[x_L,x_R]$, and $\tilde{W}(x)>0$.
  \item There is a unique $x_0(z)\in(x_L,x_R)$ such that $\tilde{D}(x_0(z))=0$.
  The point $(x_0(z),y(x_0(z),z),z)\in B$ defines a full-equality matrix in
  $\mathcal{F}_{\le}$ with fifth pivot $g(z):=\tilde{W}(x_0(z))=2H_{11}$.
  \item Every matrix $A\in\mathcal{F}_{<}$ with $(x,y,z)\in B$ satisfies
  $y=y(x,z)$, $x<x_0(z)$, and $0 < p_5(A) = \tilde{W}(x) < g(z)$.
\end{enumerate}
\end{lemma}

\begin{proof}
The exact rational interval certificate (\texttt{pivot/region.py}) establishes
$W>0$, $P_x>0$, and $P_y>0$ throughout $B$, and verifies the following
face enclosures for every $z\in[z_L,z_R]$:
\begin{itemize}
  \item At $x=x_L=-0.6177$: $P(x_L,-0.778957,z)<0$,
  $P(x_L,-0.778955,z)>0$, and $D(x_L,y,z)<0$ for all
  $y\in[-0.778957,-0.778955]$.
  \item At $x=x_R=-0.6174$: $P(x_R,-0.779307,z)<0$,
  $P(x_R,-0.779304,z)>0$, and $D(x_R,y,z)>0$ for all
  $y\in[-0.779307,-0.779304]$.
\end{itemize}
Fix $z\in[z_L,z_R]$ and any $x\in[x_L,x_R]$. Because $P_x>0$ and $P_y>0$ on
$B$ and $y_L=-0.7794 < -0.779307$ while $y_R=-0.7789 > -0.778955$,
monotonicity in $x$ and $y$ gives
\begin{align*}
  P(x,y_L,z) &\le P(x_R,y_L,z) < P(x_R,-0.779307,z) < 0, \\
  P(x,y_R,z) &\ge P(x_L,y_R,z) > P(x_L,-0.778955,z) > 0.
\end{align*}
By the intermediate value theorem and $P_y>0$, there exists a unique
$y(x,z)\in(y_L,y_R)$ solving $P(x,y(x,z),z)=0$. By the implicit function
theorem, $y(x,z)$ is $C^\infty$ with $\partial y/\partial x=-P_x/P_y$.

Differentiating along $y(x,z)$ yields
$\tilde{D}'(x) = D_x - D_y P_x/P_y$ and $\tilde{W}'(x) = W_x - W_y P_x/P_y$.
Exact rational interval arithmetic encloses these derivatives on all of $B$ in
\[
  D_x - D_y\frac{P_x}{P_y} \in (1.85,\, 1.88),
  \qquad
  W_x - W_y\frac{P_x}{P_y} \in (0.76,\, 0.78),
\]
proving Part~2. At $x=x_L$ the root $y(x_L,z)$ lies in $(-0.778957,-0.778955)$
where $D<0$, so $\tilde{D}(x_L)<0$; at $x=x_R$ the root $y(x_R,z)$ lies in
$(-0.779307,-0.779304)$ where $D>0$, so $\tilde{D}(x_R)>0$. Strict
monotonicity of $\tilde{D}$ gives a unique zero $x_0(z)\in(x_L,x_R)$. At
$(x_0(z),y(x_0(z),z),z)$, stage~4 satisfies $H_{11}>0$, $H_{12}=H_{11}$,
$H_{21}=-H_{11}$, and $H_{22}=H_{11}$, so all ties and dominance bounds hold
at all stages (by Lemma~\ref{lem:reconstruct}), proving Part~3.

Finally, any matrix $A\in\mathcal{F}_{<}$ with $(x,y,z)\in B$ satisfies $P=0$,
hence $y=y(x,z)$, and satisfies $r=H_{22}/H_{11}<1$. Because $H_{11}>0$, this
means $\tilde{D}(x)=D(x,y(x,z),z)<0=\tilde{D}(x_0(z))$, so $x<x_0(z)$ and
$0<p_5(A)=\tilde{W}(x)<\tilde{W}(x_0(z))=g(z)$.
\end{proof}

\begin{proposition}[Exact bound at full equality]\label{prop:comparator}
For every $z\in[z_L,z_R]$, the comparator pivot satisfies $g(z)\le g_*$.
\end{proposition}

\begin{proof}
At a comparator $(x_0(z),y_0(z),z)\in B$, write $g=g(z)$, $s=-G_{12}=za$,
$w=G_{31}$, and $c=G_{33}=2-(x+z)v$. From the stage-4 Schur formulas
\[
  H_{21} = -q + \frac{sw}{q},\qquad H_{22} = c + w,
\]
and the full-equality relations $H_{11}=H_{22}=-H_{21}=g/2$, we obtain
$w=g/2-c$ and hence $(q+s)(g/2)=q^2+sc$. Direct simplification gives
$q+s = 2+z-z^2$ and $q^2+sc = 4+2z(1-x-2z)$. Since $z>0$ on $B$, this uniquely
determines
\begin{equation}\label{eq:recover-x}
  x = 1 - 2z + \frac{2}{z} - \frac{g(2+z-z^2)}{4z}.
\end{equation}
Next, $H_{11}=1+s+yz-kd$ is affine in $y$ for fixed $(x,z)$ because
$k=(q-za-y(z+v))/(e-d)$ is affine in $y$ while $s,d,e,q,a,v$ do not depend on
$y$. Writing $H_{11}=\alpha(x,z)+\beta(x,z)y$, the interval certificate
(\texttt{pivot/comparator.py}) verifies $\beta>0$ throughout $B$, so
$H_{11}=g/2$ uniquely recovers $y=(g/2-\alpha(x,z))/\beta(x,z)$ as a rational
function of $(g,z)$ upon substituting~\eqref{eq:recover-x}.

Substituting these rational expressions for $x$ and $y$ into the remaining
identity $w = 1 + z^2/y + (b-q)/k = g/2 - c$ and multiplying by
$-1024\,z^4(z-1)^2 aky$ (every factor of which is well-defined and finite at
the comparator because $z\in[z_L,z_R]\subset(0,1)$ and $a,k,-y,\beta,e-d>0$
on $B$) yields an integer polynomial equation $E(g,z)=0$. The checker
(\texttt{pivot/comparator.py}, archived in \texttt{pivot/comparator.json})
derives $E$ directly from these Schur identities, verifies
$\deg_z E = 10$ and $\deg_g E = 7$, checks that the leading coefficient in $z$
is
\begin{equation}\label{eq:leading-z}
  \operatorname{lc}_z E(g) = (g-8)(g-6)(g-4)^2(g-2)(g^2-12g+40),
\end{equation}
and computes the exact polynomial discriminant in $z$:
\begin{equation}\label{eq:discriminant}
  \operatorname{Disc}_z E(g)
  = -295147905179352825856\,(g-6)(g-4)^{32}(g-2)^6(g^2-12g+40)P_5(g),
\end{equation}
matching~\cite[eq.~(2.14)]{CEU2026}.

Let $g_U := 4.14$ and consider the compact rectangle
$R := [g_L,g_U]\times[z_L,z_R]$, where $g_L$ is the rational lower bound
in~\eqref{eq:gstar-bracket}. The interval certificate establishes
$4.13 < W < g_U$ throughout $B$, so $g(z)<g_U$ for every $z\in[z_L,z_R]$.
Exact M\"obius substitutions and Descartes' rule (checked in both Python and
Lean~4) certify:
\begin{itemize}
  \item $E(g,z_L)$ and $E(g,z_R)$ each have no root for $g\in[g_L,g_U]$;
  \item $\operatorname{lc}_z E(g)$ has no root for $g\in[g_L,g_U]$;
  \item $E(g_U,z)$ has no root for $z\in[z_L,z_R]$ (independently cross-checked
  by a Sturm count).
\end{itemize}

Suppose for contradiction that $g(\hat z) > g_*$ for some
$\hat z\in[z_L,z_R]$. Then $(g(\hat z),\hat z)$ lies in the compact zero set
$Z := \{(g,z)\in R : E(g,z)=0\}$, so $Z$ is nonempty and contains a point
$(g_{\max},z_{\max})$ maximizing $g$. Necessarily
$g_* < g(\hat z) \le g_{\max} \le g_U$. By the face exclusions above,
$g_{\max}<g_U$ and $z_L < z_{\max} < z_R$. Moreover, because $P_5$ has a
unique root $g_*$ on $[4,5]$ and the remaining factors in~\eqref{eq:discriminant}
have no root on $(4,5)$ (note $g^2-12g+40=(g-6)^2+4>0$), we have
$\operatorname{Disc}_z E(g_{\max})\ne 0$. Since
$\operatorname{lc}_z E(g_{\max})\ne 0$ as well, the univariate polynomial
$z\mapsto E(g_{\max},z)$ has exact degree $10$ and no multiple roots, so
$\partial E/\partial z(g_{\max},z_{\max})\ne 0$. By the implicit function
theorem, the zero locus of $E$ continues smoothly as $z=z(g)$ to values
$g > g_{\max}$ inside $R$, contradicting the maximality of $g_{\max}$.
Hence $g(z)\le g_*$ for all $z\in[z_L,z_R]$.
\end{proof}

Combining Lemma~\ref{lem:branch}(4) and Proposition~\ref{prop:comparator}
proves Part~1 of Theorem~\ref{thm:box}: every $A\in\mathcal{F}_{<}$ in $B$
satisfies $0 < p_5(A) < g(z) \le g_*$.

\begin{remark}\label{rem:typo}
In the rendered TeX for equation~(2.13) of~\cite{CEU2026}, the $z^0$ row of
the large factor contains an extra $+97792g$ preceding $+229376g$ relative to
the derived eliminant $E(g,z)$ archived in \texttt{pivot/comparator.json}. Our
proof uses only the polynomial $E(g,z)$ derived directly from the Schur
identities and verifies its discriminant~\eqref{eq:discriminant} directly
against $P_5$.
\end{remark}

\section{Sharpness and Attainment}\label{sec:sharpness}

We now prove Parts~2 and~3 of Theorem~\ref{thm:box}.

\begin{proof}[Proof of Parts~2 and~3 of Theorem~\ref{thm:box}]
For each $z\in[z_L,z_R]$, Lemma~\ref{lem:branch} gives a unique point
$(x_0(z),y_0(z))\in(x_L,x_R)\times(y_L,y_R)$ where $P=0$ and $D=0$. The
Jacobian determinant of $(D,P)$ with respect to $(x,y)$ is
\[
  \det\begin{pmatrix} D_x & D_y \\ P_x & P_y \end{pmatrix}
  = P_y\left(D_x - D_y\frac{P_x}{P_y}\right) > 0
\]
throughout $B$. Since $P$ and $D$ are $C^\infty$ on an open neighborhood of
$B$ and $(x_0(z),y_0(z))$ is strictly interior in $(x,y)$, the implicit
function theorem and uniqueness in $B$ imply that $z\mapsto (x_0(z),y_0(z))$
is $C^\infty$ on $[z_L,z_R]$. Consequently, the comparator pivot
$g(z)=2H_{11}(x_0(z),y_0(z),z)$ is continuous on $[z_L,z_R]$ and $C^1$ on
$(z_L,z_R)$, satisfies $4.13 < g(z) < 4.14$ (from the box enclosure of $W$),
and satisfies $E(g(z),z)=0$ for all $z\in[z_L,z_R]$.

Set $c := 4.13251707863247283$ and $z_M := 0.45322491 \in (z_L,z_R)$. Exact
Sturm counts (checked in \path{verification/pivot_attainment.py}) and exact
Descartes sign variations (checked in Lean~4 by \texttt{attainment\_fibers})
certify that for each $z\in\{z_L,z_M,z_R\}$, the univariate polynomial
$g\mapsto E(g,z)$ is nonzero at $4.13$, $c$, and $4.14$ and has a unique real
root in $(4.13,4.14)$, with:
\begin{itemize}
  \item the unique root at $z=z_L=0.45322490$ lying in $(4.13,c)$;
  \item the unique root at $z=z_R=0.45322501$ lying in $(4.13,c)$;
  \item the unique root at $z=z_M=0.45322491$ lying in $(c,4.14)$.
\end{itemize}
Therefore $g(z_L) < c < g(z_M)$ and $g(z_R) < c < g(z_M)$.

By continuity on $[z_L,z_R]$, $g(z)$ attains its maximum at some interior
point $z_0\in(z_L,z_R)$. At $z_0$ we have $g'(z_0)=0$, so differentiating
$E(g(z),z)=0$ gives
\[
  \frac{\partial E}{\partial g}(g(z_0),z_0)\,g'(z_0)
  + \frac{\partial E}{\partial z}(g(z_0),z_0) = 0
  \implies \frac{\partial E}{\partial z}(g(z_0),z_0) = 0.
\]
Thus $z_0$ is a multiple root of $z\mapsto E(g(z_0),z)$. Because
$\operatorname{lc}_z E(g)$ in~\eqref{eq:leading-z} has no root in
$(4.13,4.14)$, the discriminant~\eqref{eq:discriminant} must vanish at
$g(z_0)$. Its unique root in $(4.13,4.14)$ is $g_*$, so $g(z_0)=g_*$. The
corresponding matrix $A_*\in\mathcal{F}_{\le}$ at $(x_0(z_0),y_0(z_0),z_0)\in B$
satisfies all pivot inequalities and ties (including $r=1$) and attains
$p_5(A_*)=g_*$, proving Part~3.

For Part~2, fix $z=z_0$ and consider the points $(x,y(x,z_0),z_0)\in B$ with
$x < x_0(z_0)$ approaching $x_0(z_0)$ from below. By Lemma~\ref{lem:branch},
these points lie in $B$, satisfy $P=0$ (hence $H_{21}=-H_{11}$) and $D<0$
(hence $r=H_{22}/H_{11}<1$), and satisfy all stage-1 through stage-3 dominance
inequalities together with $H_{12}=H_{11}$ and $H_{11}>0$. By continuity,
$r\to 1$ as $x\uparrow x_0(z_0)$, so $-1\le r<1$ (and hence stage-4 dominance)
holds for all $x$ sufficiently close to $x_0(z_0)$. Because $W(x,y(x,z_0),z_0)$
is continuous and strictly increasing in $x$, we have $p_5(x)\uparrow g(z_0)=g_*$,
establishing $\sup\{p_5(A): A\in\mathcal{F}_{<},\, (x,y,z)\in B\}=g_*$.
\end{proof}

\section{Necessary Condition Over the Full Relaxed Family}\label{sec:global}

We now prove Theorem~\ref{thm:global} directly from the stage entries, without
assuming any of the local reconstruction denominators $1+x$, $y$, $e-d$, or $k$
are nonzero.

\begin{proof}[Proof of Theorem~\ref{thm:global}]
Let $A\in\mathcal{F}_{<}$. Stage-1 normalization gives $x,z\in[-1,1]$ and
$|a_{13}|\le 1$. From the prescribed entries~\eqref{eq:template},
$F_{41}=-1-z$, so the tie $F_{41}=-F_{11}$ yields $F_{11}=b:=1+z$ and hence
$a_{22}=F_{11}+x=1+x+z$. Stage-1 dominance $|a_{22}|\le 1$ gives
$x+z\in[-2,0]$. Since $b=F_{11}$ is a nonzero pivot and $z\in[-1,1]$, we have
$b>0$ (so $z>-1$).

Using only division by $F_{11}=b>0$, the ties $F_{23}=-b$, $F_{21}=b$, and
$F_{41}=-b$ yield $a_{14}=-z$, $F_{13}=-1+xz$, $F_{43}=-1+z^2$, and hence
\[
  G_{12} = F_{23} - F_{13} = -z(1+x),
  \qquad
  G_{32} = F_{43} + F_{13} = -2 + z(x+z).
\]
The tie $G_{32}=-G_{11}$ gives $G_{11}=q:=2-z(x+z)$. If $z\ge 0$, then
$x+z\le 0$ implies $z(x+z)\le 0$ and thus $q\ge 2>0$; if $-1<z<0$, then
$x+z\ge -2$ implies $z(x+z) = (-z)(-(x+z)) < 1\cdot 2 = 2$ and thus $q>0$.
Hence $q>0$ in all cases.

Set $w := G_{31}$. Since $G_{32}=-q$, the stage-4 Schur update gives
$H_{21}=-q - wG_{12}/q$, and the tie $H_{21}=-H_{11}$ gives
\[
  H_{11} = q + \frac{wG_{12}}{q}.
\]
Stage-3 dominance yields $|w|\le q$ and $|G_{12}|\le q$, so
$|wG_{12}/q|\le q$ and $H_{11}\ge 0$. Because $H_{11}$ is a nonzero pivot,
$H_{11}>0$. Using $H_{12}=H_{11}$, $H_{21}=-H_{11}$, and $H_{22}=rH_{11}$, the
fifth pivot is
\[
  p_5 = H_{22} - \frac{H_{21}H_{12}}{H_{11}} = (1+r)H_{11}.
\]
Since $r\ge -1$, $H_{11}>0$, and $p_5\ne 0$, we have $p_5>0$ (so $|p_5|=p_5$).

We now bound $H_{11}$ in two cases:
\begin{itemize}
  \item \emph{Case $z\ge 0$.} Because $1+x\ge 0$ and $z\ge 0$, we have
  $|G_{12}| = z(1+x)$. Since $|w|/q\le 1$,
  \[
    H_{11} \le q + |G_{12}| = \bigl(2 - z(x+z)\bigr) + z(1+x) = 2 + z(1-z).
  \]
  At $z=1$ or $z=0$, this gives $H_{11}\le 2$ and hence
  $p_5=(1+r)H_{11} < 2\cdot 2 = 4 < g_*$.

  \item \emph{Case $-1<z\le 0$.} Since $z\le 0$ and $x+z\le 0$, we have
  $z(x+z)\ge 0$, so $q\le 2$. Moreover $G_{12}=-z(1+x)\ge 0$, so stage-3
  dominance gives $0\le G_{12}/q \le 1$. From $F_{21}=b$ and $F_{41}=-b$, the
  stage-3 updates give $q = G_{11} = F_{22}-F_{12}$ and
  $w = G_{31} = F_{42}+F_{12} = 1 - za_{13} + F_{12}$. Combining
  $F_{22}\le b$ (stage-2 dominance) and $|a_{13}|\le 1$ (stage-1 dominance)
  with $|z|=-z$ gives
  \[
    w = 1 - za_{13} + F_{22} - q \le 1 + |z| + b - q = 1 - z + (1+z) - q = 2 - q.
  \]
  Since $G_{12}/q\ge 0$, $w\le 2-q$, and $2-q\ge 0$ with $G_{12}/q\le 1$, we
  obtain
  \[
    H_{11} = q + w\frac{G_{12}}{q} \le q + (2-q)\frac{G_{12}}{q} \le q + (2-q) = 2.
  \]
  Consequently $p_5 = (1+r)H_{11} < 4 < g_*$.
\end{itemize}

Therefore any matrix $A\in\mathcal{F}_{<}$ with $p_5(A)>g_*$ (note $g_*>4$)
must have $0<z<1$ and satisfy
\[
  g_* < p_5 = (1+r)H_{11} \le (1+r)\bigl(2+z(1-z)\bigr),
\]
establishing~\eqref{eq:necessary}. Finally, since $1+r>0$ and
$2+z(1-z)\le 9/4$ on $(0,1)$,~\eqref{eq:necessary} implies
$g_* < \frac{9}{4}(1+r)$, i.e., $r > 4g_*/9 - 1 \approx 0.83667$; and since
$1+r<2$ and $2+z(1-z)>0$,~\eqref{eq:necessary} also implies
$2+z(1-z)>g_*/2$, i.e., $z(1-z) > g_*/2 - 2 \approx 0.06626$.
\end{proof}

\section{Computational Verification and Formalization}\label{sec:repro}

All algebraic identities, exact rational interval enclosures, symbolic
eliminations, discriminant factorizations, and real-root isolations are
reproduced by running \texttt{python3 -B replay.py} (Python~3.12,
SymPy~1.14.0) from the companion repository. The replay script checks SHA-256
hashes of the frozen problem specification, benchmark coefficients, and
verifier (\path{verification/FROZEN.sha256}), regenerates and compares all
archived JSON certificates (\texttt{baseline.json}, \texttt{witness.json},
\texttt{region.json}, \texttt{comparator.json}), runs the witness sub-box
check (\path{verification/pivot_subbox.py}) and attainment root-count
check (\path{verification/pivot_attainment.py}), and verifies that the
integer polynomial inputs in the Lean formalization match the archived
benchmark and derived eliminant (\path{verification/check_lean_inputs.py}).

In addition, Lean~4.34.1 formally checks the two files
\path{formalization/RatioBound.lean} and
\path{formalization/BenchmarkSigns.lean} with zero \texttt{sorry}s and only
standard axioms, verifying:
\begin{itemize}
  \item the rational order-arithmetic deductions behind~\eqref{eq:necessary}
  and the strict comparator inequality $p_5 < g(z) \le g_*$;
  \item all exact integer M\"obius-transformed sign variations and endpoint
  signs for the benchmark brackets $[4,5]$ and $[g_L,g_H]$, the four
  comparator rectangle face and leading-coefficient exclusions, and the
  attainment fiber crossings at $z_L$, $z_M$, and $z_R$.
\end{itemize}
The symbolic derivation of the eliminant $E$ and its discriminant, Descartes'
rule bridging sign variations to real roots, and the real-variable calculus
arguments (implicit function theorem and compactness) are verified in Python
and the proofs above.

\begin{thebibliography}{2}

\bibitem{CEU2026}
J.~Chen, A.~Edelman, and J.~C.~Urschel,
IMA Journal of Numerical Analysis (2026),
\href{https://doi.org/10.1093/imanum/drag039}{doi:10.1093/imanum/drag039},
\href{https://arxiv.org/abs/2602.20390}{arXiv:2602.20390v2}.

\bibitem{CEUCode}
J.~Chen, A.~Edelman, and J.~C.~Urschel,
Companion exact computation repository (\texttt{growth-5x5}, pinned commit
\texttt{bd6d13af}),
\href{https://github.com/jamiecjx/growth-5x5/tree/bd6d13af571d0c454178dac24e753f7cac7a3727}{\path{github.com/jamiecjx/growth-5x5}}.

\end{thebibliography}

\end{document}
